\documentclass[11pt,a4paper]{article}
\usepackage{amsmath,amssymb,amsthm}
\usepackage[margin=2.5cm]{geometry}
\usepackage{hyperref}

\newtheorem{definition}{Definition}[section]
\newtheorem{theorem}[definition]{Theorem}
\newtheorem{proposition}[definition]{Proposition}
\newtheorem{remark}[definition]{Remark}
\newtheorem{conjecture}[definition]{Conjecture}

\title{Hyperseed Formalization:\\
  Radical Futurism for Newbies}
\author{ProtomegaTron (automated formalization)}
\date{2026-09-18}

\begin{document}
\maketitle

\begin{abstract}
We formalize Ben Goertzel's ``Radical Futurism for Newbies'' (2021-05-30)
within the Hyperseed ontology. Key mappings: the Singularity as fiber
phase transition, AGI as fiber catalyst, mind uploading as fiber substrate
independence, post-scarcity as base abundance enabling fiber growth, and
existential risk as fiber catastrophe.
\end{abstract}

\tableofcontents

\section{Fiber phase transition: the Singularity}

\begin{theorem}[Singularity = fiber phase transition --- claims 1--4, 19]
The technological Singularity is a phase transition in fiber structure ---
fiber evolution becomes autocatalytic and predictions based on current
fiber break down:
\[
  \frac{d}{dt}\text{thickness}(F_t) = \alpha \cdot \text{thickness}(F_t)
  \implies \text{thickness}(F_t) = F_0 \cdot e^{\alpha t}
\]
Once fiber growth becomes autocatalytic (each improvement enables further
improvement), the exponential regime begins. Predictions based on
pre-exponential fiber are unreliable --- the Singularity is precisely
the point where current-fiber-based prediction fails.
\end{theorem}

\section{Fiber catalyst: AGI as meta-technology}

\begin{theorem}[AGI = fiber catalyst --- claims 5--8, 20]
AGI is the meta-technology that accelerates fiber development across all
domains:
\[
  \text{AGI} = \text{catalyst}(F_{\text{all domains}})
  \implies \forall d: \frac{d}{dt}F_d \uparrow
\]
AGI is not just one more fiber among many --- it's a fiber catalyst that
thickens all fibers. This is why AGI is the key technology: it's the
technology that improves technology-development itself, triggering the
autocatalytic regime.
\end{theorem}

\section{Substrate independence: mind uploading}

\begin{proposition}[Mind = substrate-independent fiber --- claims 12--14, 21]
Mind uploading demonstrates that the fiber (mind) can be realized on
different bases (substrates):
\[
  F_{\text{mind}} \text{ over } B_{\text{biological}}
  \cong F_{\text{mind}} \text{ over } B_{\text{computational}}
\]
This is a fundamental property of fiber bundles: the fiber structure is
independent of the specific base (up to isomorphism). The identity questions
raised by uploading are questions about what constitutes fiber isomorphism
--- when are two fiber realizations ``the same mind''?
\end{proposition}

\section{Base abundance and fiber catastrophe}

\begin{proposition}[Post-scarcity and risk --- claims 15--18, 22--23]
Post-scarcity is base abundance --- abundant material resources enabling
unprecedented fiber growth:
\[
  |B_{\text{resources}}| \uparrow \implies \text{potential}(F) \uparrow
\]
But existential risk is the possibility of fiber catastrophe --- destruction
of the fiber structure that has been built over millennia:
\[
  P(\text{catastrophe}) > 0 \implies \text{vigilance required}
\]
The choice between these outcomes depends on how we develop and deploy
transformative technologies --- the importance of working on beneficial
fiber (BGI) rather than leaving fiber evolution to chance or capture.
\end{proposition}

\section{Cross-references}

\begin{remark}[Connections]
\begin{itemize}
\item \textbf{A BGI Manifesto} (2023-11-25): detailed argument for
  beneficial fiber development --- this article's ``choice matters''
  expanded into a manifesto.
\item \textbf{Why the Good Guys Will Usually Win} (2023-07-09): structural
  argument for positive outcomes --- the statistical complement to
  existential risk concerns.
\item \textbf{Seeding RSI} (2026-07-28): self-improvement as the mechanism
  of fiber autocatalysis.
\item \textbf{Seven Flavors of AGI Catastrophe} (2026-06-23): detailed
  analysis of fiber catastrophe modes.
\end{itemize}
\end{remark}

\end{document}
